\section{Evaluation Protocol and Hyperparameter Defence}
\label{sec_meth_evaluation}

This section defines the evaluation protocol used for every result in
Chapter~\ref{chap_4}
and records the hyperparameter-selection procedure used for the PG-GAT experiments.
The same core evaluation logic is retained across three settings:
synthetic evaluation,
COCO evaluation with annotated keypoints,
and COCO end-to-end evaluation with detected keypoints.
The data distribution and source of the keypoints change between these settings,
but the grouping read-out,
label-alignment procedure,
and PGA definition remain fixed.
This common protocol makes the reported results directly interpretable
as changes in grouping performance
rather than changes in the evaluation procedure.

Hyperparameter selection is treated as part of the experimental
methodology
rather than as an informal implementation detail.
The principal PG-GAT architectural,
optimisation,
loss,
and fine-tuning choices are supported by Bayesian sweeps
over explicitly defined search spaces,
with the selection metric and stopping procedure fixed for each sweep.
The resulting configurations are then evaluated under the same protocol
as the corresponding baselines and ablations.
This separates the search used to select a configuration
from the evaluations used to characterise its behaviour.

Reproducibility is maintained through experiment logging -
every result reported in Chapter~\ref{chap_4} is associated
with a recorded run and its complete configuration,
including model hyperparameters,
optimiser settings,
dataset split,
checkpoint-selection criterion,
and evaluation parameters.
Consequently,
the numerical comparisons in the results chapter can be traced
to concrete experimental runs
rather than to unlogged manual configurations.
The subsections below specify the dataset-specific evaluation regimes
and the Bayesian sweep protocol used to select the final PG-GAT configuration.

\subsection{Datasets and splits}
\label{sec_meth_evaluation_splits}

Four evaluation views are reported across three data regimes,
all using Pose Grouping Accuracy (PGA) from Equation~\eqref{eq:pga}
as the grouping metric.

\textbf{Synthetic validation:}
The synthetic validation set consists of the 100 scenes defined in
Section~\ref{sec_meth_data_synth}.
Synthetic-validation PGA is used as the model-selection criterion
during synthetic training -
the checkpoint with the highest validation PGA is retained
for subsequent evaluation.
The same metric is used as the optimisation target for the Bayesian
architecture and loss sweeps described below.

\textbf{Synthetic test:}
The held-out synthetic test set contains a further 100 scenes
that are not used for optimisation or checkpoint selection.
Synthetic-test PGA is reported alongside validation performance
to measure generalisation beyond the split used for model selection.

\textbf{COCO val with ground-truth keypoints:}
The oracle-keypoint evaluation uses the COCO~\cite{lin2014coco}
\texttt{val2017} person-keypoint annotations
rather than detections from an upstream pose estimator.
Images are first filtered to contain at least one person
with labelled keypoints,
yielding $2{,}346$ images.
Of these,
$2{,}307$ contain sufficient valid keypoints to form a scoreable
grouping instance,
and the ground-truth-keypoint evaluations are reported over these
$2{,}307$ scenes.

The retained set contains both single-person and multi-person images.
Single-person scenes are valid inputs,
but require no non-trivial identity separation
and therefore raise the absolute PGA.
For transparency,
Chapter~\ref{chap_4} additionally reports a stricter diagnostic
restricted to \texttt{min\_people}~$\geq2$,
containing $1{,}248$ multi-person images.
The full filtered set remains the headline COCO evaluation
because the corresponding HigherHRNet reference results
are evaluated on the full distribution -
the multi-person subset is reported alongside it to expose performance
when an actual grouping decision is required.

\textbf{COCO val end-to-end:}
The end-to-end evaluation replaces annotated keypoints
with detections from the upstream pose estimator.
Candidate images are selected using a stricter pycocotools filter
requiring at least one person with at least three visible keypoints.
This retains $2{,}258$ of the $5{,}000$ \texttt{val2017} images.
An image contributes to a detector-specific PGA result only when
at least two detections can subsequently be matched to ground-truth
keypoints under the protocol of
Section~\ref{sec_meth_integration_e2e_metric}.
Because the detections differ between models,
the final number of scoreable scenes is detector-dependent:
$2{,}165$ for HigherHRNet,
$2{,}196$ for YOLO11x-pose,
and $2{,}121$ for RTMO-L.

For the primary end-to-end experiment,
the keypoints are produced by HigherHRNet-W32-512 as described in
Section~\ref{sec_meth_integration_detection}.
The same detected keypoints are grouped once by HigherHRNet's native
associative-embedding mechanism
and once by PG-GAT,
allowing their difference to be attributed to the grouping stage
rather than to detection recall.
The detector-swap evaluations repeat the same protocol with
YOLO11x-pose and RTMO-L as described in
Section~\ref{sec_meth_integration_swap}.
These end-to-end experiments are the evaluations that directly exercise
the detector-independent interface of PG-GAT.

\subsection{Run identification and the W\&B logging layer}
\label{sec_meth_evaluation_wandb}

Weights \& Biases (W\&B)~\cite{wandb} is used as the experiment-tracking
and provenance layer for the results reported in Chapter~\ref{chap_4}.
For each training run,
W\&B records the model configuration,
training and validation metrics,
source-code state,
and associated output artifacts under a persistent run identifier.
Every run cited in Chapter~\ref{chap_4} belongs to the
\texttt{multiperson-pose-grouping} project
and is assigned an eight-character run identifier.
The same identifier is used as the corresponding local output-directory
identifier in the \texttt{outputs/<name>/<run\_id>/} hierarchy
and is recorded in the experimental evidence table referenced
from Chapter~\ref{chap_4}.
This provides a common key linking a reported result to its dashboard
record,
local outputs,
configuration,
and checkpoint.

Three properties of this logging scheme are important for
reproducibility.
First,
whenever a run produces a checkpoint that improves its validation
criterion,
the checkpoint is uploaded as a W\&B artifact.
Successive versions are tracked through the auto-updated
\texttt{latest} alias,
while a post-hoc promotion script (\texttt{promote\_best.py})
assigns the project-level \texttt{best} alias after comparing
the relevant runs.
Named aliases including \texttt{champion-arch},
\texttt{champion-loss},
\texttt{champion-hyperbolic},
and \texttt{meng-headline} identify the canonical artifacts
associated with the principal architecture,
loss,
hyperbolic,
and final-model results.
These aliases provide stable names for the selected checkpoints,
while the underlying artifact versions retain their full provenance.

Second,
the complete run configuration is serialised when the W\&B run
is initialised.
Architectural hyperparameters,
optimiser settings,
dataset options,
random seeds,
and experiment-specific switches are therefore recoverable
from the experiment record independently of the local output directory.
This makes the logged run,
rather than a later verbal reconstruction,
the source of truth for determining which configuration produced
a reported result or whether a particular setting was evaluated.

Third,
every run used to support a numerical result in
Chapter~\ref{chap_4} is tagged \texttt{dissertation-final},
together with a tag identifying its experimental role,
such as \texttt{ablation},
\texttt{sweep},
or \texttt{finetune}.
The common tag provides a single queryable subset of the W\&B project
containing the experimental evidence used by the dissertation.
Consequently,
each reported result can be traced from the chapter table
to a run identifier,
from that identifier to its complete configuration and metric history,
and from the run to the exact checkpoint artifact used for evaluation.

\subsection{Bayesian sweep methodology}
\label{sec_meth_evaluation_sweeps}

Each principal PG-GAT hyperparameter underlying the headline model is
supported by a sweep over a predefined operating range.
Four sweeps are used across Section~\ref{sec_meth_pggat_training} and
Section~\ref{sec_meth_hyperbolic_hypotheses}:
an architecture sweep,
a loss-and-regularisation sweep,
a hyperbolic-architecture sweep,
and a COCO fine-tuning sweep.
Although their search spaces differ,
the four sweeps follow the same experimental protocol.

\textbf{Search strategy:}
All four sweeps use Bayesian optimisation through the W\&B sweep
agent.
Hyperband early termination~\cite{li2018hyperband} is used where
applicable to stop configurations that remain uncompetitive after the
specified number of validation iterations.
At the available budgets of $12$--$32$ runs per sweep,
Bayesian optimisation provides a sample-efficient alternative
to exhaustive grid search,
while early termination limits the compute spent on configurations
that are unlikely to become competitive.
This combination keeps each search within the compute budget available
on the experimental hardware.

\textbf{Selection metric:}
The architecture,
loss-and-regularisation,
and hyperbolic sweeps are ranked by synthetic-validation PGA
from Equation~\eqref{eq:pga}.
This is also the metric used during synthetic training
for checkpoint selection,
so the criterion used to rank sweep configurations is consistent
with the criterion used to retain the best checkpoint within each run.

The fine-tuning sweep applies the same principle after the
synthetic-pre-trained architecture has been fixed -
only the fine-tuning variables are allowed to change,
while the architecture selected by the preceding sweep remains constant.
The separation between architecture selection and fine-tuning is deliberate.
It prevents improvements caused by adaptation to the COCO distribution
from being attributed retrospectively to an architectural choice.

COCO oracle-keypoint and end-to-end results are computed
in a separate evaluation pass on the selected checkpoints.
They are therefore reported as downstream evaluations
rather than additional optimisation objectives for the architecture,
loss,
and hyperbolic searches.

\textbf{Search-space sizing:}
The bounds of every searched variable are fixed in the corresponding
W\&B sweep configuration and reproduced with the sweep results in
Chapter~\ref{chap_4}.
The ranges are chosen to cover the practically relevant operating region,
from compact configurations comparable to those used in prior
GNN-on-pose work
to the largest configurations permitted by the batch-size-$4$
memory budget of the experimental hardware.
Values outside these predefined ranges are not evaluated -
the reported optimum should therefore be interpreted as the best
configuration found within the stated search space
rather than as a global optimum over all possible PG-GAT architectures.

\textbf{Post-sweep evaluation:}
After each sweep,
the selected configuration is associated with a canonical
\texttt{champion-<sweep>} artifact alias.
The corresponding checkpoint is then passed through the common
evaluation pipeline to obtain the downstream COCO metrics required
by Chapter~\ref{chap_4}.
The resulting provenance chain is
\begin{equation}
    \text{sweep run}
    \;\longrightarrow\;
    \text{checkpoint artifact}
    \;\longrightarrow\;
    \texttt{champion}~\text{alias}
    \;\longrightarrow\;
    \text{reported evaluation}.
    \label{eq:sweep_provenance}
\end{equation}
Each stage is represented in the W\&B experiment record and,
where applicable,
in the corresponding local output directory.

\subsection{The headline provenance}
\label{sec_meth_evaluation_provenance}

The headline PG-GAT checkpoint is the product of a sequential
model-selection procedure
rather than a single isolated training run.
The architecture sweep first selects the number of message-passing
layers,
attention-head count,
hidden dimension,
output dimension,
joint-type embedding dimension,
and dropout.
With that architecture fixed,
the COCO fine-tuning sweep searches the fine-tuning duration,
learning rate,
and weight decay.
The resulting model therefore has a two-stage provenance -
architecture selection followed by fine-tuning selection.

The loss-and-regularisation sweep provides complementary evidence
for the contrastive margin and repulsion-head configuration
used by the model,
while the hyperbolic sweep performs the corresponding architecture
search for the hyperbolic variant.
These searches do not alter the provenance
of the headline Euclidean checkpoint -
they test whether important alternative settings within their respective
design axes produce a stronger configuration.

The canonical artifact produced by the final fine-tuning stage is
identified by the W\&B alias \texttt{meng-headline},
and its source run is identified by \texttt{b5noyg7t}.
The results reported for the headline PG-GAT model in
Section~\ref{sec_results_headline} are derived from this checkpoint
and can therefore be traced from the reported metric to a specific
artifact,
run configuration,
and training history.

This provenance is the basis of the hyperparameter-defensibility claim
made in this work.
The principal architectural and optimisation choices that determine
the headline PG-GAT checkpoint are either selected or validated
through explicit logged searches over stated ranges
rather than being justified retrospectively from the final test result.